% EncoderModule
\section{【応用】EncoderModule — 割り込み連携パターン}

\begin{patternbox}{EncoderModule で学べるパターン}
\begin{itemize}
    \item ハードウェア割り込み（ISR）とvolatile変数の連携
    \item IRAM\_ATTR による割り込みハンドラのRAM配置
    \item ISR→volatile変数→updateInput()→SystemDataのデータフロー
    \item deinit()による割り込みの解除
\end{itemize}
\end{patternbox}

\subsection{割り込みハンドラの設計}

\begin{lstlisting}[style=cppstyle, caption={EncoderModule.h — 割り込み関連}]
class EncoderModule : public IModule {
private:
    EncoderConfig _config;
    ModuleTimer   _sampleTimer;

    // ISRが更新する変数（volatile必須）
    volatile int32_t _isrCount = 0;

    // updateInput()内で使う前回値
    int32_t _prevCount = 0;

    // ISRはstatic関数でなければattachInterruptに渡せない
    static void IRAM_ATTR _isrHandlerA();
    static EncoderModule* _instance;

    void IRAM_ATTR _onPulseA();
    // ...
};
\end{lstlisting}

\begin{notebox}[なぜstaticが必要なのか？ — 初学者向け解説]
C++のクラスのメンバ関数（例: \texttt{\_onPulseA()}）は、内部的に「どのオブジェクトの関数か」という情報（thisポインタ）を持っています。しかし、\texttt{attachInterrupt()}が受け付けるのはC言語形式の「普通の関数ポインタ」だけで、thisポインタを渡す仕組みがありません。

そこで以下の2段構えで解決します。

\begin{enumerate}
    \item \textbf{static関数}を用意する — static関数はthisポインタを必要としないので、割り込みに登録できる
    \item \textbf{staticなインスタンスポインタ}を用意する — static関数の中から「どのオブジェクトか」を知るための橋渡し
\end{enumerate}

\begin{lstlisting}
// 1. staticポインタ: 「どのエンコーダか」を覚えておく
static EncoderModule* _instance;

// 2. static関数: 割り込みに登録できる
static void _isrHandlerA() {
    _instance->_onPulseA();  // ここで実際の処理に橋渡し
}

// init()で登録
_instance = this;
attachInterrupt(pin, _isrHandlerA, RISING);
\end{lstlisting}

これは「C++のオブジェクト指向とC言語の割り込みの橋渡し」パターンです。割り込みを使わない通常のモジュールでは不要なので、初めて見る場合はこのパターンだけ覚えれば十分です。
\end{notebox}

\begin{cautionbox}[ISR（割り込みサービスルーチン）の制約]
\begin{itemize}
    \item \texttt{attachInterrupt()}に渡す関数は\textbf{static関数}でなければならない
    \item ISR内では\texttt{Serial.print()}やメモリ確保等の重い処理は禁止
    \item ESP32等では\texttt{IRAM\_ATTR}を付けてRAMに配置する必要がある（ボード依存）
    \item ISRから参照する変数は\texttt{volatile}を付ける
\end{itemize}
\end{cautionbox}

\subsection{ISR→volatile→updateInput()の流れ}

\begin{lstlisting}[style=cppstyle, caption={EncoderModule.cpp — ISRとupdate()}]
// ISRハンドラ（IRAM_ATTR = RAM配置、高速アクセス）
void IRAM_ATTR EncoderModule::_isrHandlerA() {
    if (_instance) {
        _instance->_onPulseA();
    }
}

void IRAM_ATTR EncoderModule::_onPulseA() {
    // A相の立ち上がり時にB相を読み取って方向を判定
    if (digitalRead(_config.pinB) == HIGH) {
        _isrCount++;   // CW
    } else {
        _isrCount--;   // CCW
    }
}

// updateInput()は入力フェーズで呼ばれる
void EncoderModule::updateInput(SystemData& data) {
    // volatile変数を一度だけ読み取り
    int32_t currentCount = _isrCount;

    data.encoder.count = currentCount;

    // 速度計算（一定間隔で実行）
    if (_sampleTimer.getNowTime() >= _config.sampleIntervalMs) {
        unsigned long elapsed = _sampleTimer.getNowTime();
        _sampleTimer.setTime();

        int32_t delta = currentCount - _prevCount;
        _prevCount = currentCount;

        float revolutions = (float)delta / _config.pulsesPerRev;
        float seconds = (float)elapsed / 1000.0f;
        data.encoder.rpm = (seconds > 0)
            ? (revolutions / seconds * 60.0f) : 0.0f;
        data.encoder.isValid = true;
    }
}
\end{lstlisting}

\begin{pointbox}[volatile変数を一度だけ読み取る]
\texttt{int32\_t currentCount = \_isrCount;}で値をローカル変数にコピーしてから使う。
volatile変数を関数内で複数回参照すると、参照のたびに値が変わる可能性がある。
一度だけ読み取ることで、同一updateInput()内での値の一貫性を保証する。
\end{pointbox}

\begin{pointbox}[BleModuleとの共通パターン]
BleModule（コールバック）とEncoderModule（ISR）は、
「メインタスク外からのデータ受け渡し」という同じ問題を解決している。
どちらも\textbf{volatile変数（またはフラグ）}を介して安全にデータを渡す。
BLEはバッファ+フラグ、エンコーダはカウンタ変数という違いはあるが、原理は同じ。
\end{pointbox}

\subsection{deinit()で割り込みを解除}

\begin{lstlisting}[style=cppstyle, caption={EncoderModule.cpp — deinit()}]
void EncoderModule::deinit() {
    detachInterrupt(digitalPinToInterrupt(_config.pinA));
    _instance = nullptr;
    Serial.println("[Encoder] deinit OK");
}
\end{lstlisting}

\begin{pointbox}[割り込み解除の重要性]
割り込みを解除しないままスリープに入ると、
スリープ中にも割り込みが発生してウェイクアップの原因になる可能性がある。
また、\texttt{\_instance = nullptr}にすることで、
deinit後にISRが呼ばれてもnullチェックで安全にスキップされる。
\end{pointbox}

\begin{notebox}[複数インスタンスへの拡張]
現在の実装はstaticポインタ1つ（\texttt{\_instance}）のため、
エンコーダモジュールは1つしか使えない。
複数使う場合は\texttt{static EncoderModule* \_instances[MAX]}の配列に拡張し、
チャネルIDでインデックスする方式に変更する。
\end{notebox}

% ─── API仕様 ───────────────────────────────
\subsection{API仕様}


\begin{apibox}{EncoderModule --- ロータリーエンコーダ入力モジュール}
\textbf{定義ファイル:} \texttt{lib/EncoderModule/EncoderModule.h / .cpp}\\
\textbf{分類:} 入力モジュール（\texttt{updateInput()}のみオーバーライド）\\
\textbf{対象デバイス:} インクリメンタル型ロータリーエンコーダ（2相出力）

ハードウェア割り込みでA相パルスをカウントし、B相の状態から回転方向を判定する。
周期的にRPM（回転速度）を計算して\texttt{SystemData}に反映する。
\end{apibox}

\subsubsection{機能概要}

\begin{itemize}
    \item ハードウェア割り込み（A相）によるパルスカウント
    \item B相によるCW/CCW方向判定
    \item 周期的RPM計算（\texttt{ModuleTimer}ベース）
    \item 累積パルス数の保持（正:CW、負:CCW）
    \item \texttt{deinit()}による割り込み解除
\end{itemize}

% --- Config ---
\subsubsection{Config構造体}

\begin{funcblock}{EncoderConfig}
エンコーダの接続ピンと計測パラメータを定義する。

\begin{longtable}{l l l p{4.8cm}}
\toprule
\textbf{メンバ} & \textbf{型} & \textbf{制約・既定値} & \textbf{説明} \\
\midrule
\endhead
\texttt{pinA}             & \texttt{uint8\_t}       & ---       & エンコーダA相ピン（割り込み対応ピンを使用すること） \\
\texttt{pinB}             & \texttt{uint8\_t}       & ---       & エンコーダB相ピン \\
\texttt{pulsesPerRev}     & \texttt{uint16\_t}      & 例: 20    & 1回転あたりのパルス数 \\
\texttt{sampleIntervalMs} & \texttt{unsigned long}  & 例: 50    & RPM計算間隔 [ms] \\
\bottomrule
\end{longtable}
\end{funcblock}

% --- Data ---
\subsubsection{Data構造体}

\begin{funcblock}{EncoderData}
エンコーダの計測結果を保持する。

\begin{longtable}{l l l p{5.0cm}}
\toprule
\textbf{メンバ} & \textbf{型} & \textbf{初期値} & \textbf{説明} \\
\midrule
\endhead
\texttt{count}   & \texttt{int32\_t} & \texttt{0}     & 累積パルス数。正:CW（時計回り）、負:CCW（反時計回り） \\
\texttt{rpm}     & \texttt{float}    & \texttt{0.0f}  & 回転速度 [rpm] \\
\texttt{isValid} & \texttt{bool}     & \texttt{false} & RPM計算が1回以上行われたら\texttt{true} \\
\bottomrule
\end{longtable}
\end{funcblock}

% --- メソッド ---
\subsubsection{メソッド}

\begin{funcblock}{EncoderModule(const EncoderConfig\& config)}
\begin{tabularx}{\textwidth}{l X}
\toprule
\textbf{項目} & \textbf{内容} \\
\midrule
引数   & \texttt{config} --- \texttt{EncoderConfig}へのconst参照 \\
動作   & Configを\texttt{\_config}にコピー保持する \\
\bottomrule
\end{tabularx}
\end{funcblock}

\begin{funcblock}{bool init()}
\begin{tabularx}{\textwidth}{l X}
\toprule
\textbf{項目} & \textbf{内容} \\
\midrule
戻り値   & \texttt{bool} --- 常に\texttt{true} \\
動作     & A相・B相ピンを\texttt{INPUT\_PULLUP}で設定。A相にRISINGエッジの割り込みハンドラを登録。シングルインスタンス参照（\texttt{\_instance}）を自身に設定 \\
\bottomrule
\end{tabularx}
\end{funcblock}

\begin{funcblock}{void updateInput(SystemData\& data)}
\begin{tabularx}{\textwidth}{l X}
\toprule
\textbf{項目} & \textbf{内容} \\
\midrule
書き込み先     & \texttt{data.encoder} \\
タイミング制御 & \texttt{ModuleTimer}で\texttt{\_config.sampleIntervalMs}ごとにRPMを計算 \\
動作           & ISRが更新した\texttt{volatile}カウンタを読み取り、前回値との差分からRPMを計算。\texttt{count}・\texttt{rpm}・\texttt{isValid}を更新 \\
\bottomrule
\end{tabularx}
\end{funcblock}

\begin{funcblock}{void deinit()}
\begin{tabularx}{\textwidth}{l X}
\toprule
\textbf{項目} & \textbf{内容} \\
\midrule
動作   & A相ピンの割り込みを\texttt{detachInterrupt()}で解除する \\
\bottomrule
\end{tabularx}
\end{funcblock}

\begin{notebox}[RPM計算式]
RPM計算は以下の式で行う:

\vspace{0.3em}
$$RPM = \frac{\Delta count}{pulsesPerRev} \times \frac{60000}{sampleIntervalMs}$$
\vspace{0.3em}

\noindent ここで $\Delta count$ はサンプリング間隔内のパルス数変化量である。
\end{notebox}

\begin{cautionbox}[シングルインスタンス制約]
\texttt{attachInterrupt()}にstatic関数を渡すため、内部で\texttt{static EncoderModule* \_instance}を使用している。
同時に使用できるEncoderModuleのインスタンスは\textbf{1つのみ}である。
複数エンコーダが必要な場合はISRディスパッチの仕組みを拡張する必要がある。
\end{cautionbox}

% --- 使用例 ---
\subsubsection{使用例}

\begin{lstlisting}[style=cppstyle, caption={EncoderModuleの使用例}]
// ProjectConfig.h
const EncoderConfig ENCODER_CONFIG = {
    .pinA             = 39,  // GPIO39: A相
    .pinB             = 40,  // GPIO40: B相
    .pulsesPerRev     = 20,  // 1回転20パルス
    .sampleIntervalMs = 50,  // 50msごとに速度計算
};

// ロジックフェーズでの使用
void applyPattern(SystemData& data) {
    if (data.encoder.isValid) {
        float speed = data.encoder.rpm;
        // 回転速度に応じた制御
    }
}
\end{lstlisting}
